\documentclass[12pt]{article}
\usepackage[margin=1in]{geometry}
\usepackage{amsmath,amssymb,amsthm}
\usepackage{natbib}
\usepackage{hyperref}
\usepackage{enumitem}

\title{Summary of ``Institution Building without Commitment'' \\
\large Bassetto, Huo, and R\'ios-Rull (\textit{American Economic Review}, 2024)}
\author{Notes prepared by Thomas J.\ Sargent and Claude}
\date{March 2026}

\begin{document}
\maketitle

\tableofcontents

\bigskip

\section{Overview}

\citet{bassetto2024institution} propose a theory of \emph{gradualism} in the implementation of
desirable policies.  The central puzzle they address is: why are many policies that experts
widely agree are beneficial---carbon taxes to address climate change, reductions in
trade barriers, and even extensions of civil rights---adopted slowly rather than
implemented all at once?

Their answer centers on environments featuring \emph{time inconsistency}: because
good policies typically impose short-run costs in exchange for long-run benefits,
the current policymaker always prefers to delegate the costly initial steps to
future policymakers.  The paper introduces a new equilibrium concept---\textbf{organizational
equilibrium}---designed to capture the gradual emergence of cooperation among a
sequence of decision-makers, each of whom faces the temptation to delay.  The
equilibrium generates policy paths that gradually transit toward an allocation
that weighs both short- and long-run concerns, stopping short of the first-best
(Ramsey) outcome but substantially improving upon the myopic predictions of
Markov equilibria.

The paper applies the theory to two major contemporary policy challenges: the
slow emergence of carbon taxes and the decades-long reduction of global tariff
rates.


\section{The Substantive Issue: Gradualism}

\subsection{The Core Tension}

Many policy environments share a common structure: implementing a good policy
today imposes immediate costs on some agents, while the benefits accrue in the
future.  Examples include:
\begin{itemize}
  \item \textbf{Climate change:} Reducing carbon emissions requires costly adjustments
    to production and energy use today, while the benefits of a cooler planet are
    enjoyed by future generations.
  \item \textbf{Trade liberalization:} Lowering tariffs improves long-run allocative
    efficiency and capital accumulation, but workers in import-competing sectors suffer
    immediate wage losses.
  \item \textbf{Civil rights, European integration:} Institution-building requires
    current sacrifices for diffuse future benefits.
\end{itemize}

When policymakers are short-sighted---either because they have shorter horizons
than the population at large, or because a utilitarian planner aggregates
heterogeneous discount rates \citep{jackson2014present,jackson2015collective}---a
\emph{time inconsistency} problem arises.  The ``quasi-geometric'' or
$(\beta,\delta)$-discounting framework \citep{strotz1956myopia, phelpspollak1968,
laibson1997golden} captures this through a discount function of the form
\[
  1,\; \delta\beta,\; \delta\beta^2,\; \delta\beta^3,\;\ldots
\]
with $\delta < 1$, so that the first period's consumption is weighted more heavily
relative to all future periods.  Each period's decision-maker would like
\emph{future} decision-makers to behave cooperatively (e.g., save more, tax carbon
more, lower tariffs), but is herself tempted to free-ride on future cooperation.

\subsection{Why Not Immediate Reform?}

Under \emph{commitment} (the Ramsey solution), the period-0 agent can impose an
entire path of future actions.  This typically yields an allocation where the
initial decision-maker ``front-loads'' the benefits of her privileged position:
she consumes more (saves less, taxes carbon less) in period~0 and then imposes a
desirable path on all successors.  This treats the initial decision-maker as
special---she gets a better deal than any successor.

The opposite extreme is a \emph{Markov equilibrium}, in which each
decision-maker plays a best response to the strategies of successors, conditioning
only on payoff-relevant state variables.  This yields a constant (typically
myopic) policy that fails to internalize any long-run benefits.

Bassetto, Huo, and R\'ios-Rull argue that \emph{neither} benchmark is
satisfactory.  The Ramsey outcome gives too much power to the initial
decision-maker.  The Markov equilibrium is too pessimistic, ignoring the
possibility that cooperation can be sustained even without commitment.  What
actually happens in practice---gradual improvement---lies between these
extremes.


\section{Organizational Equilibrium: The Formal Framework}

\subsection{The Environment}

The paper studies economies with a sequence of decision-makers (``governments,''
``selves,'' or ``generations'') indexed by time.  There is a physical state
variable $k \in K$ evolving according to $k_{t+1} = F(k_t, a_t)$, where
$a_t \in A$ is the action chosen by the period-$t$ agent.  Period-$t$
preferences are given by $U(k_t, a_t, a_{t+1}, a_{t+2}, \ldots)$.

A key technical assumption is \textbf{weak separability}: the total payoff
decomposes as
\[
  U(k_t, a_t, a_{t+1}, \ldots) = v\bigl(k_t,\; V(a_t, a_{t+1}, \ldots)\bigr),
\]
where $v$ is strictly increasing in its second argument.  The ``action payoff''
$V$ depends only on the sequence of actions, not on the state.  This structure
ensures that preferences over action sequences are independent of the
inherited state, making welfare comparisons across periods meaningful.

\subsection{The Three Requirements}

Starting from the set of subgame-perfect equilibria, an organizational
equilibrium is defined by imposing three requirements:

\begin{enumerate}
  \item \textbf{No Restarting (Symmetry):}  The action payoff
    $V(a_t, a_{t+1}, \ldots)$ must be the \emph{same} for all periods~$t$ along
    the equilibrium path, and no other symmetric equilibrium achieves a higher
    payoff.  Intuitively, no decision-maker receives a more favorable treatment
    than any other.  This rules out the Ramsey outcome, in which the initial
    agent enjoys a privileged position.  The condition extends
    \citeauthor{kocherlakota1996reconsideration}'s (\citeyear{kocherlakota1996reconsideration})
    notion of ``reconsideration proofness'' to environments with state variables,
    using weak separability to factor out the state.

  \item \textbf{Optimality:}  Among all symmetric, state-independent equilibria
    satisfying no-restarting, the organizational equilibrium achieves the highest
    possible constant action payoff.

  \item \textbf{No Delay:}  The initial agent must not prefer to ``sit out''
    one period---i.e., play her best one-shot action (the Markov action)---and
    let the equilibrium process start from the next period.  Formally:
    \[
      V(a_0^*, a_1^*, a_2^*, \ldots) \ge V(a, a_0^*, a_1^*, \ldots)
      \quad \text{for all } a \in A.
    \]
    This condition ensures that cooperation cannot be so generous at the start
    as to tempt the first agent to delay.  It is the key driver of gradualism.
\end{enumerate}

\subsection{The Alternative Game: Why Period~0 Is Not Special}

The authors provide a game-theoretic micro-foundation in which the no-delay
condition emerges not as a refinement but as a \emph{requirement for a
sequential equilibrium}.  They modify the standard game as follows: the actions
of past players are \emph{unobservable} until an ``organization'' is set up to
record past play.  The opportunity to establish such record-keeping arrives at a
stochastic, unobserved time~$\hat t$.  Agent~$\hat t$ can either establish the
organization (making the history observable to successors) or pass; any future
agent can also discontinue the organization, erasing history.

In this game, if the no-delay condition were violated, the agent who first has
the opportunity to establish record-keeping would prefer to ``pass the buck''---
pretending nothing has happened and hoping that the next agent will start the
cooperation.  Hence the no-delay condition is necessary for a sequential
equilibrium, not merely a convenient refinement.

The authors describe this as the ``organization'' interpretation: an ongoing,
uninterrupted set of agents who choose to go along with a plan rather than set
up their own organization from scratch.  Governments can ``break with the
past,'' effectively restarting history.

\subsection{Key Propositions and Characterization}

Under technical assumptions (weak separability, compactness, a recursive
structure on continuation values), the paper establishes:

\begin{enumerate}[label=(\roman*)]
  \item \textbf{Existence} (Proposition~1): An organizational equilibrium exists.

  \item \textbf{Action payoff equals the best constant-action payoff}
  (Proposition~5(i)): The organizational equilibrium achieves the same payoff as
  the \emph{best constant action} profile $V(a^*, a^*, a^*, \ldots)$.  This is the
  action that best balances short- and long-run considerations when held fixed
  forever.

  \item \textbf{Intermediate welfare} (Proposition~5(ii)): The payoff lies
  strictly between the Markov equilibrium and the Ramsey outcome (unless there
  is no time inconsistency).

  \item \textbf{Gradual convergence} (Proposition~5(iii)--(iv)): The continuation
  value $\hat V(a_t, a_{t+1}, \ldots)$ increases monotonically over time,
  converging to the best constant-action value.  Convergence is \emph{not
  immediate}: the initial action differs from the steady-state action.  The
  speed of convergence slows as the economy approaches the steady state.

  \item \textbf{Recursive computation} (Proposition~4, Corollary~1): The
  organizational equilibrium can be computed by (a) finding the best constant
  action~$a^*$ and the associated action payoff~$V^*$, then (b) solving a
  recursive difference equation backward from~$a^*$ subject to the no-delay
  condition pinning down the initial action~$a_0$.
\end{enumerate}

\subsection{Comparison with Other Equilibrium Concepts}

\begin{itemize}
  \item \textbf{Ramsey:}  Gives the initial agent commitment power; typically
    features a phase of front-loaded consumption followed by a jump to a
    permanently high saving rate.  Violates no-restarting: the initial agent
    gets a higher action payoff than successors.

  \item \textbf{Markov:}  Each agent best-responds myopically.  Satisfies
    no-restarting and no-delay trivially (constant policy), but yields the
    lowest payoff.

  \item \textbf{Best subgame-perfect (trigger strategies):}  Can approximate the
    Ramsey outcome for patient agents (folk theorem), but requires \emph{grim
    triggers}---reversion to dominated outcomes---to sustain cooperation.  Also
    violates no-restarting.

  \item \textbf{Organizational:}  Cooperation builds gradually.  The punishment
    for deviation is proportional---``setting back the clock'' rather than
    reverting to worst-case outcomes.  No grim triggers required.
\end{itemize}


\section{The Motivating Example: Quasi-Geometric Discounting Growth Model}

The paper introduces the equilibrium concept via a growth model with log utility,
Cobb-Douglas production $f(k) = k^\alpha$, full depreciation, and
$(\beta,\delta)$-discounting.

In this model, the ``action'' is the saving rate $s_t = k_{t+1}/k_t^\alpha$, and
weak separability holds:  the lifetime utility decomposes into
\[
  U(k_t, s_t, s_{t+1}, \ldots) = \frac{\alpha(1 - \alpha\beta + \delta\alpha\beta)}{1 - \alpha\beta} \log k_t + V(s_t, s_{t+1}, \ldots).
\]
Three benchmarks emerge in closed form:
\begin{align*}
  s^M &= \frac{\alpha\delta\beta}{1 - \alpha\beta + \delta\alpha\beta}
    && \text{(Markov)}, \\
  s^R &= \alpha\beta
    && \text{(long-run Ramsey)}, \\
  s^B &= \frac{\delta\alpha\beta}{(1-\beta+\delta\beta)(1-\alpha\beta) + \delta\alpha\beta}
    && \text{(best constant saving rate)},
\end{align*}
with $s^M < s^B < s^R$ when $\delta \in (0,1)$.

The organizational equilibrium converges to $s^B$.  The transition is governed by
a difference equation $s_{t+1} = q(s_t; V^*)$, starting from an initial saving
rate $s_0 = q(s^M; V^*)$ pinned down by no-delay.  The dynamics exhibit:
\begin{itemize}
  \item Monotonic increase from $s_0$ toward $s^B$;
  \item Decreasing speed of convergence as $s_t \to s^B$;
  \item A constant action payoff $V^*$ along the entire path.
\end{itemize}


\section{Application I: Climate Change Mitigation}

\subsection{Model}

The paper extends the climate-economy model of \citet{golosov2014optimal} by
introducing quasi-geometric discounting.  Production uses capital, labor, and an
energy input (coal), with a carbon externality entering through an atmospheric
carbon stock that damages productivity:
\[
  y_t = \exp(-\gamma q_t)\, k_t^\alpha\, n_t^{1-\alpha-\nu}\, e_t^\nu.
\]
Carbon accumulates through permanent and persistent components.  The
policymaker balances a short-sighted politician's bias ($\delta < 1$) against
the long-run damage from carbon emissions.

A key separability property holds: given an initial state, the payoff decomposes
additively into a component depending on the state $(k_0, q_{1,-1}, q_{2,-1})$,
a component $W$ depending on the saving rates, and a component $V$ depending on
labor allocation to the final-goods sector.  Because the carbon tax~$\Lambda_t$
is determined by the labor allocation, characterizing the organizational
equilibrium for labor suffices to characterize the equilibrium path of the
carbon tax.

\subsection{Steady-State Comparison}

In steady state, the carbon tax under the three equilibrium concepts ranks as:
\[
  \Lambda^M < \Lambda^B < \Lambda^R.
\]
The Markov carbon tax is the lowest (most myopic); the Ramsey tax fully
internalizes future damages; the organizational equilibrium tax lies between
them.

\subsection{Transition Dynamics}

The organizational equilibrium generates a \emph{gradually rising} carbon tax,
starting from a level near the Markov tax and converging to $\Lambda^B$.  This
contrasts with the Ramsey outcome, which features a low initial tax followed by
an abrupt jump, and with the Markov equilibrium, which maintains a constant low
tax.

\subsection{Historical Evidence}

The paper documents two pieces of empirical evidence consistent with the
gradualism prediction:

\begin{enumerate}
  \item \textbf{Scandinavian countries} (Denmark, Finland, Norway, Sweden), the
    first group to implement carbon taxes, exhibit a gradually increasing trend
    in (real) carbon tax rates over time.

  \item \textbf{Cross-country evidence:}  After controlling for the quantity
    of carbon emissions covered and GDP per capita, carbon tax rates are
    positively correlated with the number of years since implementation
    (coefficient $\hat\beta_1 = 1.11$, $p$-value~$= 0.017$).

  \item \textbf{China's evolving stance:}  The paper contrasts President
    Hu~Jintao's 2008 statement prioritizing economic development over
    environmental policy with President Xi~Jinping's 2020 pledge to reach peak
    carbon emissions by 2030 and carbon neutrality by 2060.  The shift from
    reluctance to ambition---still falling short of immediate maximum
    abatement---embodies the gradualist dynamic.
\end{enumerate}

More broadly, the succession of international agreements---the UN Framework
Convention on Climate Change (1992), the Kyoto Protocol (2005), the Paris
Agreement (2015), and COP28 (2023)---illustrates increasingly ambitious but still
gradual ratcheting up of emission standards.


\section{Application II: Tariff Policies and Trade Liberalization}

\subsection{Model}

The paper constructs a two-country trade model with the following ingredients:
\begin{itemize}
  \item Two symmetric countries, each producing two tradeable intermediate goods
    via capital and sector-specific labor, aggregated into a final good via a CES
    technology.
  \item Comparative advantage: country~1 is more productive in good~1, and vice
    versa.
  \item No labor mobility across sectors: a tariff shields the importing sector's
    workers from wage losses.
  \item An aggregate capital externality drives endogenous growth; a lower tariff
    raises the return on capital, encouraging capital accumulation and thus
    long-run growth.
  \item The policymaker maximizes average worker welfare (utilitarian over both
    sectors).
\end{itemize}

The core tension: lowering tariffs hurts sector-2 workers in the short run (they
face foreign competition) but benefits all workers in the long run through faster
growth.  Under commitment, the policymaker would always prefer setting higher
tariffs initially and lower tariffs later---making the problem time-inconsistent.

\subsection{Results}

\begin{itemize}
  \item The Markov equilibrium features a \emph{constant, high tariff}: the most myopic
    outcome, where policymakers fail to internalize the growth benefits of free trade.
  \item The Ramsey outcome features a high initial tariff followed by a rapid elimination
    of tariffs---exploiting the commitment power of the initial policymaker.
  \item The organizational equilibrium features a \emph{gradual decline} in tariffs,
    starting from an intermediate level and converging to a low (but positive)
    steady-state tariff rate.
\end{itemize}

For workers in the importing sector, the organizational equilibrium initially
delivers lower consumption than the Markov equilibrium (due to the lower tariff),
but eventually delivers higher consumption because the lower tariff path
stimulates greater capital accumulation and growth.

\subsection{Historical Evidence}

The paper documents the gradual decline in world average tariff rates from
approximately 16\% in 1950 to about 4\% in 2017, based on data from
\citet{antras2020}.  The institutional mechanism---successive rounds of GATT
negotiations, culminating in the formation of the WTO---mirrors the gradualist
logic of the organizational equilibrium.

The paper notes a particularly attractive property for analyzing trade
disruptions:  after a ``tremble'' in which a government unexpectedly raises
tariffs, the organizational equilibrium does \emph{not} require grim-trigger
punishment.  Instead, the gradual march toward cooperation simply resumes from
the new, worse starting point.  The authors view this as a useful framework for
analyzing episodes such as the reimposition of tariffs that did not lead to
all-out trade wars but rather to negotiations restarting to undo newly introduced
barriers.


\section{Connections to Other Equilibrium Refinements}

The organizational equilibrium is closely related to several strands of the
literature on equilibrium refinements in dynamic games:

\begin{itemize}
  \item \textbf{Reconsideration proofness} \citep{kocherlakota1996reconsideration}:
    In purely repeated games (no state variable), the no-restarting condition
    reduces to Kocherlakota's symmetry requirement.  Bassetto et~al.\ extend
    this to dynamic games with state variables by using weak separability to
    factor out the state and comparing only the action payoff~$V$.

  \item \textbf{Revision proofness} \citep{asheim1997,alessleet2014}: Under
    revision proofness, the set of credible punishments is larger because a
    deviating player must propose alternatives that benefit \emph{all} future
    players.  This leads to a much larger equilibrium set (for quasi-geometric
    discounting with linear preferences, all subgame-perfect paths better than
    Markov are revision-proof).

  \item \textbf{Recursive Pareto criterion} \citep{brendonellison2018}: Brendon
    and Ellison restrict the Ramsey planner to sequences satisfying a recursive
    Pareto criterion.  Their steady-state allocation coincides with the
    organizational equilibrium's limiting allocation, but their motivation is
    normative whereas Bassetto et~al.'s is positive.

  \item \textbf{APS value-set methods} \citep{abreu1990}: Proposition~4 uses
    continuation values as a state variable, similar to APS, but the symmetry
    constraint (all continuation equilibria share the same value) complicates
    the standard monotone iteration procedure.
\end{itemize}


\section{Broader Historical Motivation}

The conclusion of the paper invokes a broad set of historical parallels.  The
authors cite \citet{sargent2017honoring} and
\citet{hallsargent2014,hallsargent2018} on the gradual process by which the
United States acquired a reputation for honoring its debts.  They also note the
post-1970s emergence of low inflation as governments ``gradually learned how to
manage monetary policy without resorting to the anchor of a commodity standard.''
Social Security, which ``started as a narrow and limited program and only later
grew in size and scope,'' provides yet another example.

The paper suggests two additional applications for future research:
\begin{itemize}
  \item \textbf{The extension of civil rights} in the United States after the
    Civil War.
  \item \textbf{The gradual creation of the European Union}, echoing the
    Schumann Declaration of 1950 that opens the paper: ``Europe will not be made
    all at once, or according to a single plan.  It will be built through
    concrete achievements which first create a de facto solidarity.''
\end{itemize}

In each of these cases, the organizational equilibrium provides a lens through
which a good institution gradually emerges (unlike in a Markov equilibrium) but
remains imperfect (worse than Ramsey).  The equilibrium concept captures the idea
that ``Rome was not built in a day.''


\section{The Gradual Acquisition of US Fiscal Credibility}

Bassetto, Huo, and R\'ios-Rull cite \citet{sargent2017honoring} and
\citet{hallsargent2014,hallsargent2018} on the slow process by which the
United States acquired a reputation for honoring its debts.  This section
expands on that subtheme, tracing the gradualist logic through a sequence
of 19th-century episodes that together illustrate how fiscal credibility
was built---and occasionally threatened---without the benefit of
credible commitment.

\subsection{Hamilton's Founding Settlement (1790)}

The United States began its existence under a fiscal cloud.  The
Continental Congress and individual states had accumulated roughly
\$79~million in debts from the Revolutionary War---domestic loan-office
certificates, final-settlement certificates, Quartermaster and
Commissary certificates, foreign loans from France and the
Netherlands---many trading at 15--25 cents on the dollar
\citep[pp.~271--275]{Ferguson1961}.  Under the Articles of Confederation,
the national government lacked the taxing power to service these
obligations, and several states had begun their own assumption programs,
creating competing creditor constituencies attached to state governments
rather than to the union \citep[pp.~237--239]{Ferguson1961}.

Alexander Hamilton's \textit{First Report on the Public Credit} (January
1790) proposed a comprehensive restructuring: fund the Continental debt
at par, assume approximately \$25~million in state war debts, and service
the consolidated federal debt out of tariff and excise revenues
\citep{Hamilton1790}.  Three elements of Hamilton's plan are especially
relevant to the gradualism framework:

\begin{enumerate}
  \item \textbf{Rejection of discrimination.}  James Madison proposed
    paying current holders only the prevailing market price and
    distributing the remainder to original holders (soldiers, farmers,
    suppliers).  Hamilton opposed discrimination on the ground that it
    would destroy the transferability of government securities and
    thereby undermine the liquid public-debt market he sought to create.
    Paying current holders at par established a precedent:\ US federal
    securities were fully transferable claims, serviced without regard to
    the identity of the holder.  This was the \emph{initial action} in a
    long sequence---costly to taxpayers in the short run, but essential
    for establishing the credibility of federal debt.

  \item \textbf{The sinking fund as commitment device.}  Hamilton
    complemented assumption with a sinking fund (Act of August~12, 1790)
    governed by a board of commissioners including the Vice President,
    the Chief Justice, and the Secretary of the Treasury.  By earmarking
    surplus revenues for debt retirement and placing the fund under
    high-profile governance, Hamilton made it politically costly for
    future Congresses to divert debt-service revenues. In the language
    of organizational equilibrium, this was an institutional arrangement
    designed to raise the cost of ``restarting''---i.e., reneging on the
    cooperative path.

  \item \textbf{Assumption as equilibrium selection.}  Hamilton's plan
    was not merely a generous bailout; it was a strategic move to win a
    competition for creditor allegiance against state governments that
    had been pursuing their own assumption programs under the Articles
    of Confederation \citep{Ferguson1961, Rothbard2019}.  The Compromise
    of 1790---exchanging votes for assumption for the location of the
    national capital on the Potomac---illustrates the political
    side-payments required to shift from a state-dominant to a
    federal-dominant fiscal equilibrium.
\end{enumerate}

Hamilton's settlement was not the Ramsey outcome: he could not commit
future Congresses to maintain the sinking fund, nor could he prevent the
political temptation to inflate away the debt.  But it was far better than
the Markov equilibrium that had prevailed under the Confederation, where
each state and the national government acted myopically, and public debt
traded at deep discounts for want of any credible commitment to service
it.

\subsection{Gallatin and Jeffersonian Validation (1801--1812)}

A critical test of Hamilton's institutional design came when political power
transferred to his opponents.  Thomas Jefferson's Treasury Secretary,
Albert Gallatin, entered office intending to dismantle much of Hamilton's
fiscal-military state---yet he preserved and even strengthened the
commitment to debt service.  Gallatin systematically applied surplus
revenues to retire the federal debt, reducing it from \$83~million in 1801
to \$45~million by 1812 \citep{balack2012gallatin}.  The fact that a
Jeffersonian administration---ideologically hostile to Hamilton's
program---nonetheless honored the debt and used surpluses to retire it
was a powerful signal.  In the organizational-equilibrium framework, this
corresponds to the second and third movers along the equilibrium path
choosing to continue the cooperative sequence rather than ``restarting''
with a deviation.

\subsection{The War of 1812 and Subsequent Retirement}

The War of 1812 forced the US government back into heavy borrowing, much
of it on unfavorable terms: Congress had allowed the First Bank of the
United States' charter to expire in 1811, weakening the institutional
infrastructure for debt management.  Borrowing costs rose sharply,
especially after British victories in 1814 \citep{hallsargent2014}.

Yet the postwar period demonstrated the resilience of the gradualist
path.  Under successive administrations, the federal government applied
surplus tariff revenues to systematic debt retirement, achieving the
remarkable feat of paying off the entire national debt by January~1835.
Andrew Jackson's Treasury Secretary announced that every dollar of the
federal debt had been extinguished---the only time a major nation has
achieved this.  Whatever one's assessment of the wisdom of eliminating
the national debt entirely (Hamilton would not have approved, since he
valued a moderate funded debt as a source of financial-market liquidity),
the episode powerfully reinforced the credibility of the US government's
commitment to honor its obligations.

\subsection{The 1840s State Debt Crisis: A Clarifying Test}

The state debt crisis of the early 1840s was a critical juncture
\citep{Ratchford1941, McGrane1935}.  Nine states defaulted on debts
incurred during the 1830s infrastructure boom, and several repudiated
portions of their debt entirely (Arkansas, Florida, Michigan,
Mississippi).  European bondholders, invoking the 1790 precedent,
demanded that the federal government assume the defaulted state debts.

Congress refused.  The refusal was costly in the short run: European
creditors temporarily downgraded all American government bonds, raising
federal borrowing costs.  But it was immensely consequential for the
long-run structure of American fiscal credibility:

\begin{itemize}
  \item It clarified that the 1790 assumption had been a one-time
    constitutional settlement, not a permanent insurance scheme.
  \item It  established a credible no-bailout norm, avoiding the
    moral-hazard problem that \citet{KarekenWallace1978} later
    formalized.
  \item It triggered state-level constitutional reforms---balanced-budget
    requirements, debt ceilings, voter-approval provisions for new
    borrowing---that persist to this day \citep{Wallis2005}.
  \item Over subsequent years, creditors learned to distinguish
    sharply between federal bonds (which had never defaulted) and state
    bonds (which carried real default risk).
\end{itemize}

In the organizational-equilibrium framework, Congress's refusal
illustrates how a deviation (state defaults) does not trigger a
grim-trigger reversion to the worst outcome; instead, the federal
government's own cooperative path continues, while the defaulting states
are ``set back''---they must rebuild credibility through costly
constitutional constraints on their own borrowing.

\subsection{Civil War Finance: Statutory vs.\ Reputational Commitment}

The Civil War pushed the federal government into the most drastic fiscal
measures since the Revolution: the Legal Tender Act of February~25, 1862,
authorized \$150~million in fiat ``greenbacks'' and the issuance of 5--20
bonds (redeemable after five years, maturing in twenty).  A revealing
feature of this legislation, documented by \citet{hallsargent2014}, is
the \emph{hybrid commitment structure} that sustained borrowing:

\begin{itemize}
  \item \textbf{Statutory commitment for interest:} Section~5 of the Legal
    Tender Act required import duties to be paid in coin and earmarked
    that coin first for interest payments on bonds and notes.
  \item \textbf{Reputational commitment for principal:} The statute was
    entirely silent on whether bond \emph{principal} would be repaid in
    coin or greenbacks.  Treasury Secretary Salmon P.~Chase bridged this
    gap through marketing: his bond circulars and Jay Cooke's network of
    2,500 agents sold the 5--20s on the understanding that principal
    would be repaid in gold.
\end{itemize}

This hybrid structure is a vivid instance of gradualism under imperfect
commitment.  The government could not formally commit to gold repayment
of principal---the political coalition for the Legal Tender Act would not
have held---but it could build a \emph{reputational} expectation through
the actions of the current Treasury Secretary.  When the postwar Congress
debated the ``Ohio Idea'' (paying off 5--20 bonds in greenbacks rather
than gold), the reputational commitment was tested.  The Public Credit
Act of 1869, pledging the government to ``payment in coin or its
equivalent'' of all US obligations, converted Chase's extra-statutory
promise into durable statutory form---a further step along the
organizational-equilibrium path.

\subsection{Resumption of Specie Payments (1879)}

The Resumption Act of 1875, which committed the Treasury to redeeming
greenbacks in gold beginning January~1,~1879, completed the postwar
restoration of fiscal credibility.  The fourteen-year gap between the
suspension of convertibility (1861) and the legislated return to gold
(1875), followed by a further four-year implementation lag, is itself an
illustration of gradualism.  Each Congress along the way faced the
temptation to delay---resumption imposed real costs on debtors and
inflationists---but the succession of administrations (Johnson, Grant,
Hayes) ultimately moved, slowly and imperfectly, toward the hard-money
equilibrium.

\subsection{Summary: Gradualism as a Running Theme}

The 19th-century American fiscal experience can be read as a sequence of
steps along an organizational-equilibrium path:

\begin{center}
\begin{tabular}{lp{10cm}}
  \textbf{1790} & Hamilton's funding and assumption: the initial action
    $a_0$ that started the cooperative sequence, at a level below the
    Ramsey ideal (no commitment of future Congresses) but far above
    the Confederation's Markov equilibrium. \\[4pt]
  \textbf{1801--12} & Jeffersonian validation: successors continue the
    cooperative path rather than restarting. \\[4pt]
  \textbf{1815--35} & Postwar debt retirement: the saving rate gradually
    rises, reinforcing credibility. \\[4pt]
  \textbf{1842} & State debt crisis and federal refusal to assume:
    proportional punishment (states are set back) without grim-trigger
    reversion; no-bailout norm clarified. \\[4pt]
  \textbf{1862--65} & Civil War: hybrid statutory/reputational commitment
    under fiscal stress. \\[4pt]
  \textbf{1869} & Public Credit Act: reputational commitment converted to
    statutory form. \\[4pt]
  \textbf{1875--79} & Resumption: gradualist return to specie
    convertibility. \\[4pt]
\end{tabular}
\end{center}

No single episode corresponds to a ``big bang'' in which the United States
leapt to the Ramsey outcome.  Nor did the process proceed smoothly: the
Civil War was a massive shock that temporarily set back the cooperative
path.  But the organizational-equilibrium logic captures the essential
pattern: each successive decision-maker inherited a partially built
institution, found it (weakly) in her interest to continue rather than
restart, and left the institution slightly stronger for the next
decision-maker.  The limiting allocation---a government that could borrow
on terms comparable to Britain's, the gold standard of sovereign
credibility---was approached gradually and imperfectly, exactly as
Bassetto, Huo, and R\'ios-Rull's theory predicts.


\section{Assessment and Relation to Fiscal History}

\subsection{What the Paper Achieves}

\begin{enumerate}
  \item It offers a \emph{positive theory} of gradualism that is grounded in rigorous
    game theory but amenable to computation: the equilibrium can be found by
    solving a recursive difference equation, much as competitive equilibria are
    computed in standard macroeconomics.

  \item It dispenses with \emph{grim triggers}: punishments for deviation are
    proportional (``set back the clock''), making the equilibrium robust to
    trembles and stochastic shocks.

  \item It delivers \emph{clear comparative statics}: unlike the folk-theorem
    multiplicity of subgame-perfect equilibria, the organizational equilibrium
    provides a sharp prediction.

  \item The two quantitative applications---carbon taxes and tariffs---show that
    the theory's qualitative predictions (gradual improvement, convergence to an
    intermediate steady state) align with observed data.
\end{enumerate}

\subsection{Relevance for Fiscal History}

The organizational equilibrium framework has natural applications to fiscal
history:
\begin{itemize}
  \item The slow building of sovereign creditworthiness, as documented by
    \citet{sargent2017honoring} and \citet{hallsargent2014,hallsargent2018} for
    the United States.
  \item The gradual evolution of fiscal institutions (debt limits, balanced-budget
    rules, the separation of capital and ordinary budgets studied in
    \citealt{bassettosargent2006}).
  \item The post-WWII development of international monetary institutions, where
    successive agreements (Bretton Woods, its collapse, the Plaza Accord, the
    adoption of the euro) can be read as gradual, imperfect institution-building.
\end{itemize}

The paper's key insight---that a succession of self-interested policymakers can
build institutions gradually even without commitment, but that the resulting
institutions will always fall short of the first-best---resonates deeply with
the themes of fiscal history.


\bibliographystyle{plainnat}
\bibliography{bassetto_et_al_summary}

\end{document}
